\documentclass{article}
\usepackage[utf8]{inputenc}
\usepackage[margin=1.0in]{geometry}
\usepackage{amsmath}
\usepackage{amssymb}
\usepackage{amsthm}
\usepackage{amsfonts}
\usepackage{booktabs}
\usepackage[parfill]{parskip}
\usepackage[dvipsnames]{xcolor}
\usepackage{hyperref}
\hypersetup{
    colorlinks=true,
    linkcolor=RoyalPurple,
    urlcolor=RoyalPurple,
    citecolor=RoyalPurple
}

\newtheorem{theorem}{Theorem}[section]
\newtheorem{proposition}[theorem]{Proposition}
\newtheorem{lemma}[theorem]{Lemma}
\newtheorem{corollary}[theorem]{Corollary}
\theoremstyle{definition}
\newtheorem{definition}[theorem]{Definition}
\newtheorem{assumption}[theorem]{Assumption}
\theoremstyle{remark}
\newtheorem{remark}[theorem]{Remark}

\newcommand{\FF}{\mathcal{F}}
\newcommand{\R}{\mathbb{R}}
\newcommand{\C}{\mathbb{C}}
\newcommand{\Z}{\mathbb{Z}}
\newcommand{\Ghat}{\widehat{G}}
\newcommand{\Hom}{\mathrm{Hom}}

\title{The Per-Pixel Steerable Bispectrum\\for Non-Abelian Groups}
\author{Johan Mathe}
\date{\today}

\begin{document}
\maketitle

\section{Introduction}

The $G$-bispectrum is a well-established invariant for signals on groups and homogeneous spaces. Across cyclic groups, finite non-abelian groups, compact groups, spheres, disks, and balls, the same pipeline applies: the group Fourier transform yields equivariant spectral coefficients, the triple correlation provides invariance, and the bispectrum---as the Fourier transform of the triple correlation---furnishes a generically complete invariant under the group action \cite{kakarala1992, kakarala2009, kondor2008}. The reconstruction procedure is also uniform: one recovers the trivial Fourier component from $\beta_{\rho_0,\rho_0}$, bootstraps a generator irrep via $\beta_{\rho_0,\rho_1}$, and iterates through the Kronecker product table to reach all remaining irreps from bispectral coefficients of the form $\beta_{\rho_1, \rho_k}$.

However, a fundamental obstacle arises when one attempts to apply this bispectral machinery to the output of a \textit{steerable convolutional network}. In steerable CNNs \cite{weiler2019}, the convolution filter is decomposed on the intertwiner basis $\Hom^G(\pi, \rho_G \otimes \rho_{\mathrm{in}}^*)$, and the dimension of this basis depends on the choice of output representation $\rho_G$. When $\rho_G = \rho_0$ (the trivial representation), the intertwiner space $\Hom^G(\pi, \rho_0)$ has dimension $m^\pi_{\rho_0}$---the multiplicity of the trivial irrep in the action on the kernel support. For example, $D_4$ acting on a $3 \times 3$ grid gives $L = 3$ basis functions, all in the $A_1$-sector, while the bispectrum requires Fourier blocks at all 5 irreps ($A_1, A_2, B_1, B_2, E$). The spectral data at the nontrivial irreps is simply absent. This is the \textit{trivial-output obstruction}: the number of intertwiner basis functions does not match the number of irreps, and the Fourier coefficients needed for the bispectrum cannot be extracted from the convolution output.

One might attempt to work around this obstruction by defining a modified bispectrum over the intertwiner coefficients themselves, or by projecting onto a different harmonic basis (e.g., disk harmonics). But these approaches introduce additional difficulties: the intertwiner basis is not in general steerable, the coefficients are not indexed by irreps, and the resulting invariant lacks the completeness guarantees of the classical $G$-bispectrum.

We resolve the obstruction by a different route: we change the \textit{output fiber representation} from $\rho_0$ to the regular representation $\rho_{\mathrm{reg}}$ of $G$. This is a choice about the output fiber type, not the spatial kernel representation $\pi$, which continues to describe the group action on the pixel grid. By the Peter--Weyl theorem, the regular representation decomposes as $\C[G] \cong \bigoplus_{\psi \in \Ghat} \psi \otimes \psi^*$, so the output fiber at each pixel carries \textit{all} isotypic components of $G$, each with multiplicity $d_\psi$. The fiber vector $\Theta(x) \in \C^{|G|}$ is canonically a function $\theta_x : G \to \C$, whose $G$-Fourier transform directly yields all Fourier blocks---no change of basis is required.

With all Fourier blocks accessible, the classical bispectral machinery applies per-pixel. We define the per-pixel $G$-bispectrum and establish:
\begin{enumerate}
    \item \textbf{Invariance} (general finite $G$): the per-pixel bispectrum is $G$-invariant and translation-equivariant (Theorem~\ref{thm:perpixel_bsp_inv}).
    \item \textbf{Completeness} (general finite $G$): by Kakarala's theorem, the per-pixel bispectrum is a generically complete invariant of the \textit{fiber signal} $\theta_x$, meaning $\theta_x$ can be reconstructed from the bispectrum up to $G$-action (Theorem~\ref{thm:completeness}).
    \item \textbf{Constructive recovery and selectivity} (specialized to $G = D_n$): we give an explicit reconstruction algorithm that bootstraps all irreps from a single generator $\rho_1$, following the same breadth-first strategy on the Kronecker table used for signals on groups. Only $O(|G|)$ bispectral entries per pixel suffice, versus $O(|G|^2)$ for the full bispectrum (Theorem~\ref{thm:selectivity}).
\end{enumerate}

A key architectural consequence is the \textit{decomposition of invariances}. One could alternatively define a Fourier transform over the full semi-direct product group $\Gamma = \Z^2 \ltimes G$, integrating over all translations and group elements simultaneously. This produces a global bispectral descriptor that is invariant to the full group $\Gamma$, but it discards the spatial index entirely and cannot be composed as a layer in a deep network. Our per-pixel approach instead handles translation equivariance through the convolutional structure and $G$-invariance through the bispectrum at each pixel, yielding a composable $G$-invariant layer that serves as a drop-in replacement for norm or gated nonlinearities.

We emphasize a precise scope for the completeness result: it applies to the per-pixel fiber signal $\theta_x : G \to \C$, not directly to the original image $f$. The steerable convolution maps $f$ to $\theta_x$ equivariantly, but the information content of $\theta_x$ about $f$ depends on the learned filter weights $\{w_l\}$. Establishing completeness for $f$ would require showing that the steerable layer, with appropriate filter weights, produces fiber signals whose Fourier blocks are generically nonsingular---a condition that we expect to hold for generic weights but do not prove here.

Throughout, we work over complex representation spaces. Real-valued signals (as arise from typical image data) are embedded in their complexification; all statements hold in particular for the real-valued case.

\section{Preliminaries}

\subsection{Groups and representations}

Let $G$ be a finite group with identity $e$. A \textit{representation} of $G$ is a group homomorphism $\rho : G \to \mathrm{GL}(V)$ for a finite-dimensional complex vector space $V$. The dimension of $\rho$ is $d_\rho = \dim V$. A representation is \textit{irreducible} (an \textit{irrep}) if $V$ has no proper $G$-invariant subspace. We write $\Ghat$ for a complete set of pairwise inequivalent irreps of $G$.

\begin{definition}[Regular representation]
The \textit{left regular representation} of $G$ on $\C[G] \cong \C^{|G|}$ is defined by
\[
[\rho_{\mathrm{reg}}(g)]_{h,h'} = \delta_{h, g h'}, \quad g, h, h' \in G.
\]
That is, $\rho_{\mathrm{reg}}(g)$ permutes the basis vectors $\{e_h\}_{h \in G}$ by left multiplication: $e_{h'} \mapsto e_{gh'}$.
\end{definition}

\begin{lemma}[Regular action as left translation on functions]\label{lem:reg_as_translation}
Under the identification $v \in \C[G] \leftrightarrow \theta_v : G \to \C$ with $\theta_v(h) = v_h$, the left regular action corresponds to left translation on functions:
\[
(\rho_{\mathrm{reg}}(r)\, \theta_v)(g) = \theta_v(r^{-1}g), \quad \forall\, r, g \in G.
\]
\end{lemma}
\begin{proof}
$[\rho_{\mathrm{reg}}(r)\, v]_g = \sum_{h'} [\rho_{\mathrm{reg}}(r)]_{g,h'}\, v_{h'} = \sum_{h'} \delta_{g, rh'}\, v_{h'} = v_{r^{-1}g} = \theta_v(r^{-1}g)$.
\end{proof}

\begin{theorem}[Peter--Weyl for finite groups]
The regular representation decomposes as
\[
\C[G] \cong \bigoplus_{\psi \in \Ghat} \psi \otimes \psi^*,
\]
where each irrep $\psi$ appears with multiplicity $d_\psi = \dim \psi$. Consequently, $\sum_{\psi \in \Ghat} d_\psi^2 = |G|$.
\end{theorem}

We denote the trivial representation by $\rho_0$, with $\rho_0(g) = 1$ for all $g \in G$.

\subsection{$G$-Fourier transform}\label{sec:gft}

\begin{definition}[$G$-Fourier transform]
Let $\theta : G \to \C$ be a signal on $G$. Its \textit{Fourier transform} at irrep $\psi \in \Ghat$ is the matrix
\begin{equation}\label{eq:gft}
\FF(\theta)_\psi = \sum_{g \in G} \theta(g)\, \psi(g)^\dagger \;\in\; \C^{d_\psi \times d_\psi}.
\end{equation}
\end{definition}

\begin{proposition}[Fourier inversion]\label{prop:inversion}
The signal $\theta$ is recovered from its Fourier transform by
\[
\theta(g) = \frac{1}{|G|} \sum_{\psi \in \Ghat} d_\psi \,\mathrm{Tr}\bigl(\psi(g)\, \FF(\theta)_\psi\bigr).
\]
\end{proposition}

\begin{proposition}[Equivariance under left translation]\label{prop:ft_equiv_prelim}
Define the left translate of $\theta$ by $r \in G$ as $\theta^r(g) = \theta(r^{-1}g)$. Then
\begin{equation}\label{eq:ft_equiv}
\FF(\theta^r)_\psi = \FF(\theta)_\psi \cdot \psi(r)^\dagger.
\end{equation}
\end{proposition}
\begin{proof}
\begin{align*}
\FF(\theta^r)_\psi
&= \sum_{g \in G} \theta(r^{-1}g)\, \psi(g)^\dagger \\
&= \sum_{g' \in G} \theta(g')\, \psi(rg')^\dagger \qquad [g' = r^{-1}g] \\
&= \sum_{g'} \theta(g')\, \psi(g')^\dagger\, \psi(r)^\dagger \\
&= \FF(\theta)_\psi \cdot \psi(r)^\dagger. \qedhere
\end{align*}
\end{proof}

\subsection{Clebsch--Gordan decomposition}\label{sec:cg}

The tensor product of two irreps $\psi_1, \psi_2 \in \Ghat$ decomposes into a direct sum of irreps:
\begin{equation}\label{eq:cg}
(\psi_1 \otimes \psi_2)(g) = C_{\psi_1,\psi_2} \left[\bigoplus_{\psi \in \psi_1 \otimes \psi_2} \psi(g)\right] C_{\psi_1,\psi_2}^\dagger, \quad \forall g \in G,
\end{equation}
where $C_{\psi_1,\psi_2}$ is a unitary matrix called the \textit{Clebsch--Gordan matrix}, and the direct sum runs over the irreducible summands appearing in $\psi_1 \otimes \psi_2$ \textit{with multiplicity} (counted according to a fixed ordering). The decomposition is recorded in the \textit{Kronecker table} of $G$.

\subsection{$G$-Triple correlation}\label{sec:tc_prelim}

\begin{definition}[$G$-Triple correlation]\label{def:tc}
The triple correlation of a signal $\theta : G \to \C$ is
\begin{equation}\label{eq:tc}
T(\theta)_{g_1, g_2} = \sum_{g \in G} \theta(g)\, \theta(g \cdot g_1)\, \theta(g \cdot g_2), \quad g_1, g_2 \in G.
\end{equation}
The right offsets $g \cdot g_1$, $g \cdot g_2$ ensure that $T(\theta)$ is invariant under \textit{left} translation, which is the action induced by the regular fiber (Lemma~\ref{lem:reg_as_translation}).
\end{definition}

\begin{proposition}[Invariance of the triple correlation]\label{prop:tc_inv}
The triple correlation is invariant under left translation: if $\theta'(g) = \theta(r^{-1}g)$ for some $r \in G$, then $T(\theta') = T(\theta)$.
\end{proposition}
\begin{proof}
\begin{align*}
T(\theta')_{g_1, g_2}
&= \sum_{g \in G} \theta(r^{-1}g)\, \theta(r^{-1}g \cdot g_1)\, \theta(r^{-1}g \cdot g_2) \\
&= \sum_{g' \in G} \theta(g')\, \theta(g' g_1)\, \theta(g' g_2) \qquad [g' = r^{-1}g] \\
&= T(\theta)_{g_1, g_2}. \qedhere
\end{align*}
\end{proof}

\subsection{$G$-Bispectrum}\label{sec:bsp_prelim}

\begin{definition}[$G$-Bispectrum]\label{def:bsp}
The bispectrum of a signal $\theta : G \to \C$ is defined, for each pair of irreps $\psi_1, \psi_2 \in \Ghat$, as
\begin{equation}\label{eq:bsp}
\beta(\theta)_{\psi_1, \psi_2}
= \bigl[\FF(\theta)_{\psi_1} \otimes \FF(\theta)_{\psi_2}\bigr]\,
C_{\psi_1,\psi_2}\,
\left[\bigoplus_{\psi \in \psi_1 \otimes \psi_2} \FF(\theta)_\psi^\dagger\right]
C_{\psi_1,\psi_2}^\dagger
\;\in\;\C^{d_{\psi_1}d_{\psi_2} \times d_{\psi_1}d_{\psi_2}},
\end{equation}
where the direct sum is over the irreducible summands of $\psi_1 \otimes \psi_2$ with multiplicity, in the same fixed ordering as the CG decomposition \eqref{eq:cg}. Extraction of individual Fourier blocks from the block-diagonal structure depends on this ordering.
\end{definition}

\begin{theorem}[Bispectrum as Fourier transform of the triple correlation]\label{thm:bsp_is_ft_tc}
The bispectrum equals the $(G \times G)$-Fourier transform of the triple correlation:
\[
\beta(\theta)_{\psi_1, \psi_2}
= \sum_{g_1 \in G} \sum_{g_2 \in G}
T(\theta)_{g_1, g_2}\, \psi_1(g_1)^\dagger \otimes \psi_2(g_2)^\dagger.
\]
\end{theorem}
\begin{proof}
Expanding the triple correlation and exchanging the order of summation:
\begin{align*}
&\sum_{g_1, g_2 \in G} T(\theta)_{g_1,g_2}\, [\psi_1(g_1)^\dagger \otimes \psi_2(g_2)^\dagger] \\
&= \sum_{g_1, g_2} \sum_{g \in G} \theta(g)\, \theta(gg_1)\, \theta(gg_2)\, [\psi_1(g_1)^\dagger \otimes \psi_2(g_2)^\dagger] \\
&= \sum_g \theta(g) \left[\sum_{g_1} \theta(gg_1)\, \psi_1(g_1)^\dagger\right] \otimes \left[\sum_{g_2} \theta(gg_2)\, \psi_2(g_2)^\dagger\right].
\end{align*}
Substituting $g_1' = gg_1$ and using $\psi_1(g_1)^\dagger = \psi_1(g^{-1}g_1')^\dagger = \psi_1(g_1')^\dagger \psi_1(g^{-1})^\dagger$:
\begin{align*}
&= \sum_g \theta(g) \left[\FF(\theta)_{\psi_1} \cdot \psi_1(g)  \right] \otimes \left[\FF(\theta)_{\psi_2} \cdot \psi_2(g) \right] \\
&= [\FF(\theta)_{\psi_1} \otimes \FF(\theta)_{\psi_2}] \sum_g \theta(g)\, [\psi_1(g) \otimes \psi_2(g)].
\end{align*}
Applying the Clebsch--Gordan decomposition \eqref{eq:cg}:
\begin{align*}
\sum_g \theta(g)\, [\psi_1(g) \otimes \psi_2(g)]
&= C_{\psi_1,\psi_2} \left[\bigoplus_{\psi} \sum_g \theta(g)\, \psi(g)\right] C_{\psi_1,\psi_2}^\dagger \\
&= C_{\psi_1,\psi_2} \left[\bigoplus_{\psi} \FF(\theta)_\psi^\dagger\right] C_{\psi_1,\psi_2}^\dagger.
\end{align*}
Combining yields \eqref{eq:bsp}.
\end{proof}

\begin{proposition}[Invariance of the bispectrum]\label{prop:bsp_inv_prelim}
The bispectrum is invariant under left translation: $\beta(\theta^r) = \beta(\theta)$ for all $r \in G$.
\end{proposition}
\begin{proof}
By Theorem~\ref{thm:bsp_is_ft_tc}, $\beta(\theta) = \FF_{G \times G}(T(\theta))$. By Proposition~\ref{prop:tc_inv}, $T(\theta^r) = T(\theta)$, so $\beta(\theta^r) = \beta(\theta)$.
\end{proof}

\subsection{Completeness of the $G$-bispectrum}

\begin{theorem}[Kakarala \cite{kakarala2009}]\label{thm:completeness_prelim}
Let $\theta, \tilde{\theta} : G \to \C$ with $\FF(\theta)_\psi$ nonsingular (i.e., $\det \FF(\theta)_\psi \neq 0$) for all $\psi \in \Ghat$. Then
\[
\beta(\theta) = \beta(\tilde{\theta})
\quad\Longleftrightarrow\quad
\exists\, r \in G \text{ such that } \tilde{\theta}(g) = \theta(r^{-1}g) \;\;\forall g \in G.
\]
\end{theorem}

The non-degeneracy condition $\det \FF(\theta)_\psi \neq 0$ defines a Zariski-open dense subset of $\C[G]$: the degenerate signals form a proper algebraic subvariety, so the condition holds generically.

\subsection{Steerable convolution}\label{sec:steerable_conv}

Consider the semi-direct product group $\Gamma = \Z^2 \ltimes G$, where $G$ acts on $\Z^2$ by $r \cdot x$ (e.g., by rotating the pixel grid). The group operation in $\Gamma$ is $(t_1, r_1)(t_2, r_2) = (t_1 + r_1 \cdot t_2,\; r_1 r_2)$, and $\Gamma$ acts on $\Z^2$ by $(t, r) \cdot x = r \cdot x + t$.

We consider signals $f : \Z^2 \to \C^{K_{\mathrm{in}}}$ equipped with the representation $T$ of $\Gamma$ induced by a representation $\rho_{\mathrm{in}}$ of $G$ on the input fiber $\C^{K_{\mathrm{in}}}$:
\begin{equation}\label{eq:induced_action}
[T_{(t,r)}\, f](x) = \rho_{\mathrm{in}}(r)\, f\bigl(r^{-1}(x - t)\bigr) \;\in\; \C^{K_{\mathrm{in}}}.
\end{equation}

A \textit{steerable convolution} maps $f$ to an output signal $\Theta : \Z^2 \to \C^{K_{\mathrm{out}}}$ carrying an output representation $\rho_G$ of $G$ on $\C^{K_{\mathrm{out}}}$. The filter $\Psi : \Z^2 \to \C^{K_{\mathrm{out}} \times K_{\mathrm{in}}}$ is decomposed on the intertwiner basis. Let $\pi$ be the representation of $G$ on the filter support (permuting pixel positions), and let $\{\Psi_l\}_{l=1}^L$ be a basis for the space $\Hom^G(\pi, \rho_G \otimes \rho_{\mathrm{in}}^*)$ of $G$-equivariant kernel maps. The filter is
\begin{equation}\label{eq:filter_decomp}
\Psi(x) = \sum_{l=1}^{L} w_l\, \Psi_l(x), \quad L = \dim \Hom^G(\pi, \rho_G \otimes \rho_{\mathrm{in}}^*),
\end{equation}
where $w_l \in \C$ are learned weights. The convolution output at pixel $x_0$ is:
\begin{equation}\label{eq:conv}
\Theta(x_0) = (\Psi * f)(x_0) = \sum_{x \in \Omega_s} \Psi(x_0 - x)\, f(x) \;\in\; \C^{K_{\mathrm{out}}}.
\end{equation}

The intertwiner constraint ensures equivariance: if $f' = T_{(t, r)}[f]$, then
\begin{equation}\label{eq:equiv_general}
\Theta'(r \cdot x_0 + t) = \rho_G(r)\, \Theta(x_0).
\end{equation}

\subsection{The trivial-output obstruction}\label{sec:obstruction}

With $\rho_G = \rho_0$ (trivial representation), the output fiber is one-dimensional and transforms trivially under $G$. The convolution output at each pixel is a scalar that is already $G$-invariant---but this invariant accesses only the trivial isotypic component ($A_1$) of the output representation. Unlike the regular fiber $\C[G]$, a trivial-output fiber contains only the trivial irrep and therefore cannot encode Fourier blocks for nontrivial irreps.

The $G$-bispectrum, however, requires Fourier blocks for every $\psi \in \Ghat$, including matrix-valued ones. For example, $D_4$ has 5 irreps ($A_1, A_2, B_1, B_2, E$) of dimensions $1, 1, 1, 1, 2$. A trivial-output convolution on a $3 \times 3$ kernel gives $L = \dim\Hom^G(\pi, \rho_0) = 3$ basis functions, all in the $A_1$-sector. The data at $A_2, B_1, B_2, E$ cannot be recovered.

\section{The regular-representation fiber}

\subsection{Definition}

We resolve the trivial-output obstruction of Section~\ref{sec:obstruction} in the model formulation by choosing the output fiber to be the \textit{regular representation} of $G$.

\begin{definition}[Regular-rep fiber]\label{def:regrep}
We set $\rho_G = \rho_{\mathrm{reg}}$, the left regular representation of $G$ on $\C[G] \cong \C^{|G|}$. The output fiber dimension is $K_{\mathrm{out}} = |G|$, and the basis vectors of the fiber are indexed by group elements $h \in G$.
\end{definition}

By the Peter--Weyl theorem, $\C[G] \cong \bigoplus_{\psi \in \Ghat} \psi \otimes \psi^*$. The fiber now contains \textit{all} isotypic components of $G$, each with multiplicity $d_\psi$. For $D_4$: $K_{\mathrm{out}} = 8 = 1^2 + 1^2 + 1^2 + 1^2 + 2^2$, and all five irreps ($A_1, A_2, B_1, B_2, E$) are represented.

\subsection{Resolution of the obstruction}

With $\rho_G = \rho_{\mathrm{reg}}$, the output fiber $\C^{K_{\mathrm{out}}} = \C^{|G|}$ is canonically identified with $\C[G]$, and the fiber vector at pixel $x$ is a function on $G$ (by the identification $v \leftrightarrow \theta_v$ of Lemma~\ref{lem:reg_as_translation}). The $G$-Fourier transform of this function directly yields all Fourier blocks---no additional change-of-basis is required to interpret the output fiber itself as a function on $G$ and apply the group Fourier transform to it.

\begin{remark}[Scope of completeness]\label{rem:scope}
The per-pixel bispectrum (defined below) is a complete invariant of the fiber signal $\theta_x : G \to \C$, not directly of the original image $f$. The steerable convolution maps $f$ equivariantly to $\theta_x$, but the information content of $\theta_x$ about $f$ depends on the learned filter weights. Establishing completeness for $f$ would require an additional analysis of the steerable layer itself.
\end{remark}

\subsection{Steerable convolution with regular-rep fiber}

The steerable convolution with regular-rep output takes the form of \eqref{eq:conv} with $K_{\mathrm{out}} = |G|$:
\begin{equation}\label{eq:conv_reg}
\Theta(x_0) = \sum_{x \in \Omega_s} \Psi(x_0 - x)\, f(x) \;\in\; \C^{|G|},
\end{equation}
where $\Psi(x) = \sum_{l=1}^L w_l\, \Psi_l(x)$ with $\Psi_l \in \Hom^G(\pi, \rho_{\mathrm{reg}} \otimes \rho_{\mathrm{in}}^*)$.

\begin{proposition}[Equivariance]\label{prop:equiv_reg}
Under the action of $(t, r) \in \Gamma$ on the input, the output satisfies
\[
\Theta'(r \cdot x_0 + t) = \rho_{\mathrm{reg}}(r)\, \Theta(x_0).
\]
In components: $\Theta'_g(r \cdot x) = \Theta_{r^{-1}g}(x)$ for all $g \in G$.
\end{proposition}
\begin{proof}
This is a special case of \eqref{eq:equiv_general} with $\rho_G = \rho_{\mathrm{reg}}$. The component form follows from Lemma~\ref{lem:reg_as_translation}: $[\rho_{\mathrm{reg}}(r)\, \Theta(x)]_g = \Theta_{r^{-1}g}(x)$.
\end{proof}

\begin{remark}\label{rem:left_translation}
The component form $\Theta'_g(r \cdot x) = \Theta_{r^{-1}g}(x)$ means that the group action on the fiber is left translation of the function $h \mapsto \Theta_h(x)$ on $G$. This is precisely the group action assumed by the $G$-Fourier transform equivariance (Proposition~\ref{prop:ft_equiv_prelim}) and the $G$-bispectrum completeness theorem (Theorem~\ref{thm:completeness_prelim}).
\end{remark}

\section{Per-pixel $G$-Fourier transform}

\subsection{Definition}

\begin{definition}[Per-pixel signal on $G$]\label{def:perpixel_signal}
At each pixel $x \in \Z^2$, define the function $\theta_x : G \to \C$ by
\begin{equation}\label{eq:theta_x}
\theta_x(h) := \Theta_h(x).
\end{equation}
\end{definition}

\begin{definition}[Per-pixel $G$-Fourier transform]\label{def:perpixel_ft}
The per-pixel $G$-Fourier transform at irrep $\psi \in \Ghat$ is
\begin{equation}\label{eq:perpixel_ft}
\widehat{\Theta}(x, \psi) := \FF(\theta_x)_\psi = \sum_{h \in G} \Theta_h(x)\, \psi(h)^\dagger \;\in\; \C^{d_\psi \times d_\psi}.
\end{equation}
\end{definition}

\begin{remark}[Relationship to the full-group Fourier transform]
Definition~\ref{def:perpixel_ft} applies the standard $G$-Fourier transform \eqref{eq:gft} to the fiber vector at each pixel. It is \textbf{not} the full-group Fourier transform over $\Gamma = \Z^2 \ltimes G$, which would integrate over all translations and group elements simultaneously and produce a global descriptor without a spatial index. In our construction, spatial translation equivariance is handled by the convolutional structure (Section~\ref{sec:steerable_conv}); the per-pixel $G$-Fourier transform handles only the $G$-action on the fiber.
\end{remark}

\begin{remark}[No information loss]
The map $\Theta(x) \mapsto \{\widehat{\Theta}(x, \psi)\}_{\psi \in \Ghat}$ is invertible by Fourier inversion on $G$ (Proposition~\ref{prop:inversion}). The total number of entries across all Fourier blocks is $\sum_{\psi} d_\psi^2 = |G|$, matching the fiber dimension.
\end{remark}

\subsection{Equivariance}

\begin{proposition}[Equivariance of per-pixel Fourier blocks]\label{prop:perpixel_ft_equiv}
Under the action of $r \in G$ on the input:
\begin{equation}\label{eq:perpixel_ft_equiv}
\widehat{\Theta}'(r \cdot x, \psi) = \widehat{\Theta}(x, \psi) \cdot \psi(r)^\dagger.
\end{equation}
\end{proposition}
\begin{proof}
By Proposition~\ref{prop:equiv_reg}, $\Theta'_g(r \cdot x) = \Theta_{r^{-1}g}(x)$. This means $\theta'_{r \cdot x}(g) = \theta_x(r^{-1}g)$, which is exactly the left translate $\theta_x^{r}$ from Proposition~\ref{prop:ft_equiv_prelim}. Therefore:
\begin{align*}
\widehat{\Theta}'(r \cdot x, \psi)
&= \FF(\theta'_{r \cdot x})_\psi \\
&= \FF(\theta_x^{r})_\psi \\
&= \FF(\theta_x)_\psi \cdot \psi(r)^\dagger \qquad\text{[by Proposition~\ref{prop:ft_equiv_prelim}]} \\
&= \widehat{\Theta}(x, \psi) \cdot \psi(r)^\dagger. \qedhere
\end{align*}
\end{proof}

\section{Per-pixel $G$-bispectrum}

\subsection{Per-pixel triple correlation}

\begin{definition}[Per-pixel $G$-triple correlation]\label{def:perpixel_tc}
At each pixel $x \in \Z^2$:
\begin{equation}\label{eq:perpixel_tc}
T(\theta_x)_{g_1, g_2} = \sum_{g \in G} \theta_x(g)\, \theta_x(g \cdot g_1)\, \theta_x(g \cdot g_2), \quad g_1, g_2 \in G.
\end{equation}
This is the $G$-triple correlation (Definition~\ref{def:tc}) applied to the per-pixel signal $\theta_x$.
\end{definition}

\begin{proposition}[Per-pixel TC invariance]\label{prop:perpixel_tc_inv}
If the input is transformed by $r \in G$, then $T(\theta'_{r \cdot x}) = T(\theta_x)$.
\end{proposition}
\begin{proof}
By Proposition~\ref{prop:equiv_reg}, $\theta'_{r \cdot x}(g) = \theta_x(r^{-1}g)$. By Proposition~\ref{prop:tc_inv} applied to $\theta_x$ with the left translation by $r$:
\begin{align*}
T(\theta'_{r \cdot x})_{g_1, g_2}
&= \sum_{g \in G} \theta_x(r^{-1}g)\, \theta_x(r^{-1}g \cdot g_1)\, \theta_x(r^{-1}g \cdot g_2) \\
&= \sum_{g' \in G} \theta_x(g')\, \theta_x(g' g_1)\, \theta_x(g' g_2) \qquad [g' = r^{-1}g]\\
&= T(\theta_x)_{g_1, g_2}. \qedhere
\end{align*}
\end{proof}

\subsection{Per-pixel bispectrum}

\begin{definition}[Per-pixel $G$-bispectrum]\label{def:perpixel_bsp}
At each pixel $x \in \Z^2$, for each pair of irreps $\psi_1, \psi_2 \in \Ghat$:
\begin{equation}\label{eq:perpixel_bsp}
\beta(x)_{\psi_1, \psi_2}
= \bigl[\widehat{\Theta}(x,\psi_1) \otimes \widehat{\Theta}(x,\psi_2)\bigr]\,
C_{\psi_1,\psi_2}\,
\left[\bigoplus_{\psi \in \psi_1 \otimes \psi_2} \widehat{\Theta}(x,\psi)^\dagger\right]
C_{\psi_1,\psi_2}^\dagger.
\end{equation}
This is the $G$-bispectrum (Definition~\ref{def:bsp}) applied to $\theta_x$, with the same multiplicity and ordering conventions.
\end{definition}

\begin{theorem}[Invariance of the per-pixel bispectrum]\label{thm:perpixel_bsp_inv}
For all $r \in G$:
\begin{equation}
\beta'(r \cdot x)_{\psi_1, \psi_2} = \beta(x)_{\psi_1, \psi_2}.
\end{equation}
That is, the per-pixel bispectrum is $G$-invariant.
\end{theorem}
\begin{proof}
By Theorem~\ref{thm:bsp_is_ft_tc}, $\beta(x)$ is the $(G \times G)$-Fourier transform of $T(\theta_x)$. By Proposition~\ref{prop:perpixel_tc_inv}, $T(\theta'_{r \cdot x}) = T(\theta_x)$. Therefore $\beta'(r \cdot x) = \beta(x)$.
\end{proof}

\begin{remark}[Translation equivariance]
The per-pixel bispectrum $\beta$ inherits translation equivariance from the convolutional structure: if the input is translated by $t \in \Z^2$, then $\beta'(x + t) = \beta(x)$. The bispectrum map $x \mapsto \beta(x)$ is therefore a translation-equivariant, $G$-invariant feature map of the fiber signal.
\end{remark}

\section{Completeness: reconstruction up to group action}

\subsection{Statement and proof}

\begin{theorem}[Completeness of the per-pixel bispectrum]\label{thm:completeness}
Let $\theta_{x_1}, \theta_{x_2} : G \to \C$ be the per-pixel fiber signals at two pixels, with $\det\bigl(\FF(\theta_{x_i})_\psi\bigr) \neq 0$ for all $\psi \in \Ghat$ and $i = 1, 2$. Then
\begin{equation}
\beta(x_1) = \beta(x_2)
\quad\Longleftrightarrow\quad
\exists\, r \in G \text{ such that } \theta_{x_2}(g) = \theta_{x_1}(r^{-1}g) \;\;\forall g \in G.
\end{equation}
Equivalently: $\Theta(x_2) = \rho_{\mathrm{reg}}(r)\, \Theta(x_1)$.
\end{theorem}
\begin{proof}
The per-pixel bispectrum $\beta(x)$ is exactly the $G$-bispectrum of $\theta_x$ (Definition~\ref{def:perpixel_bsp}). The result follows immediately from Theorem~\ref{thm:completeness_prelim} (Kakarala) applied to the functions $\theta_{x_1}$ and $\theta_{x_2}$.
\end{proof}

We now give a constructive recovery procedure specialized to $G = D_n$, which additionally yields the selectivity result of Section~\ref{sec:selectivity}.

\subsection{Irreducible representations of $D_n$}\label{sec:dn_irreps}

Let $D_n = \langle a, x \mid a^n = x^2 = e,\; xax = a^{-1} \rangle$ be the dihedral group of order $2n$, where $a$ is rotation by $2\pi/n$ and $x$ is a reflection.

\textbf{One-dimensional irreps.}
For $n$ even ($n \geq 4$), there are four 1D irreps:
\begin{alignat*}{3}
&\rho_0 : \quad &&\rho_0(a^l x^m) = 1 &&\text{(trivial),} \\
&\rho_{01} : \quad &&\rho_{01}(a^l x^m) = (-1)^m &&\text{(sign on reflections),} \\
&\rho_{02} : \quad &&\rho_{02}(a^l x^m) = (-1)^l &&\text{(sign on half-rotations),} \\
&\rho_{03} : \quad &&\rho_{03}(a^l x^m) = (-1)^{l+m}.
\end{alignat*}
For $n$ odd, only $\rho_0$ and $\rho_{01}$ exist (since $(-1)^l$ is not a well-defined character when $a^n = e$ with $n$ odd).

\textbf{Two-dimensional irreps.}
For $k = 1, \ldots, \lfloor(n-1)/2\rfloor$:
\begin{equation}\label{eq:dn_2d_irrep}
\rho_k(a^l x^m) = \begin{bmatrix} \cos\omega_k l & -\sin\omega_k l \\ \sin\omega_k l & \cos\omega_k l \end{bmatrix} \begin{bmatrix} 1 & 0 \\ 0 & -1 \end{bmatrix}^m, \quad \omega_k = \frac{2\pi k}{n}.
\end{equation}

\textbf{Dimension count.}
For $n$ even: $4 \cdot 1^2 + (n/2 - 1) \cdot 2^2 = 4 + 2n - 4 = 2n = |D_n|$.
For $n$ odd: $2 \cdot 1^2 + \frac{n-1}{2} \cdot 2^2 = 2 + 2(n-1) = 2n = |D_n|$.

\subsection{Recovering $\FF(\theta_x)_{\rho_0}$}

\begin{lemma}\label{lem:recover_rho0}
The bispectral coefficient at $(\rho_0, \rho_0)$ satisfies
\begin{equation}\label{eq:beta_00}
\beta(\theta)_{\rho_0, \rho_0} = |\FF(\theta)_{\rho_0}|^2\, \FF(\theta)_{\rho_0}.
\end{equation}
\end{lemma}
\begin{proof}
Since $\rho_0$ is the trivial representation, $\FF(\theta)_{\rho_0} = \sum_{g} \theta(g) \in \C$ is a scalar. Also $\rho_0 \otimes \rho_0 = \rho_0$, and the CG matrix is the identity. Substituting into \eqref{eq:bsp}:
\begin{align*}
\beta(\theta)_{\rho_0, \rho_0}
&= [\FF(\theta)_{\rho_0} \otimes \FF(\theta)_{\rho_0}] \cdot \FF(\theta)_{\rho_0}^\dagger \\
&= \FF(\theta)_{\rho_0}^2 \cdot \overline{\FF(\theta)_{\rho_0}} \\
&= |\FF(\theta)_{\rho_0}|^2 \cdot \FF(\theta)_{\rho_0}. \qedhere
\end{align*}
\end{proof}

\textbf{Recovery.} From $\beta(\theta)_{\rho_0,\rho_0}$:
\begin{equation}\label{eq:recover_rho0}
|\FF(\theta)_{\rho_0}| = |\beta(\theta)_{\rho_0,\rho_0}|^{1/3}, \qquad
\arg\bigl(\FF(\theta)_{\rho_0}\bigr) = \arg\bigl(\beta(\theta)_{\rho_0,\rho_0}\bigr).
\end{equation}
This uniquely determines $\FF(\theta)_{\rho_0}$.

\subsection{Recovering $\FF(\theta_x)_{\rho_1}$}

We choose the 2D irrep $\rho_1$ (defined in \eqref{eq:dn_2d_irrep} with $k = 1$) as the \textit{generator irrep}. Its tensor powers, combined via the Kronecker table, generate all irreps of $D_n$ (Proposition~\ref{prop:generation} below).

\begin{lemma}\label{lem:recover_rho1}
The bispectral coefficient at $(\rho_0, \rho_1)$ satisfies
\begin{equation}\label{eq:beta_01}
\beta(\theta)_{\rho_0, \rho_1} = \FF(\theta)_{\rho_0} \cdot \FF(\theta)_{\rho_1} \cdot \FF(\theta)_{\rho_1}^\dagger.
\end{equation}
\end{lemma}
\begin{proof}
Since $\rho_0 \otimes \rho_1 = \rho_1$ and the CG matrix for $\rho_0 \otimes \rho_1$ is the identity:
\begin{align*}
\beta(\theta)_{\rho_0, \rho_1}
&= [\FF(\theta)_{\rho_0} \otimes \FF(\theta)_{\rho_1}] \cdot \FF(\theta)_{\rho_1}^\dagger \\
&= \FF(\theta)_{\rho_0} \cdot \FF(\theta)_{\rho_1} \cdot \FF(\theta)_{\rho_1}^\dagger. \qedhere
\end{align*}
\end{proof}

\textbf{Recovery.} Since $\FF(\theta)_{\rho_0}$ is known from \eqref{eq:recover_rho0}:
\begin{equation}\label{eq:recover_rho1_gram}
\FF(\theta)_{\rho_1} \cdot \FF(\theta)_{\rho_1}^\dagger = \frac{\beta(\theta)_{\rho_0, \rho_1}}{\FF(\theta)_{\rho_0}}.
\end{equation}
The right-hand side is a known Hermitian positive semi-definite matrix (positive definite under non-degeneracy). Compute its eigenvalue decomposition $U \Lambda U^\dagger = \FF(\theta)_{\rho_1} \cdot \FF(\theta)_{\rho_1}^\dagger$, and set
\begin{equation}\label{eq:recover_rho1}
\FF(\theta)_{\rho_1} = U \Lambda^{1/2} U^\dagger \cdot W,
\end{equation}
where $W$ is unitary. This determines $\FF(\theta)_{\rho_1}$ up to a right-unitary ambiguity $W$. By Theorem~\ref{thm:completeness_prelim}, the full bispectral constraints reduce this ambiguity to a left-translation orbit of the reconstructed signal; we do not provide an independent proof of uniqueness for the constructive scheme here.

\subsection{Bootstrapping remaining irreps}

\begin{lemma}\label{lem:bootstrap}
Let $\psi_a, \psi_b \in \Ghat$ with $\FF(\theta)_{\psi_a}$ and $\FF(\theta)_{\psi_b}$ known and invertible. Then every irrep $\psi$ appearing in the tensor decomposition $\psi_a \otimes \psi_b$ can be recovered from $\beta(\theta)_{\psi_a, \psi_b}$:
\begin{equation}\label{eq:bootstrap}
\bigoplus_{\psi \in \psi_a \otimes \psi_b} \FF(\theta)_\psi^\dagger
= C_{\psi_a,\psi_b}^\dagger\,
\bigl[\FF(\theta)_{\psi_a} \otimes \FF(\theta)_{\psi_b}\bigr]^{-1}\,
\beta(\theta)_{\psi_a, \psi_b}\,
C_{\psi_a,\psi_b}.
\end{equation}
The individual blocks $\FF(\theta)_\psi^\dagger$ are read off from the block-diagonal structure, once a CG basis with fixed ordering of irreducible summands is chosen.
\end{lemma}
\begin{proof}
From \eqref{eq:bsp}, left-multiply by $[\FF(\theta)_{\psi_a} \otimes \FF(\theta)_{\psi_b}]^{-1}$ and right-multiply by $C_{\psi_a,\psi_b}$:
\begin{align*}
[\FF_{\psi_a} \otimes \FF_{\psi_b}]^{-1} \beta_{\psi_a,\psi_b}\, C_{\psi_a,\psi_b}
&= C_{\psi_a,\psi_b} \left[\bigoplus_\psi \FF_\psi^\dagger\right] C_{\psi_a,\psi_b}^\dagger\, C_{\psi_a,\psi_b} \\
&= C_{\psi_a,\psi_b} \left[\bigoplus_\psi \FF_\psi^\dagger\right].
\end{align*}
Left-multiplying by $C_{\psi_a,\psi_b}^\dagger$ (which is unitary) yields \eqref{eq:bootstrap}.
\end{proof}

\subsection{Generating all irreps for $D_n$}

We show that the generator irrep $\rho_1$ iteratively produces all irreps of $D_n$ through the bootstrap procedure.

\begin{proposition}\label{prop:generation}
For $D_n$ with $n > 2$:
\begin{enumerate}
    \item For $k = 2, \ldots, \lfloor(n-1)/2\rfloor$: the 2D irrep $\rho_k$ appears in $\rho_1 \otimes \rho_{k-1}$.
    \item The 1D irrep $\rho_{01}$ appears in $\rho_1 \otimes \rho_1$.
    \item For $n$ even: the 1D irreps $\rho_{02}, \rho_{03}$ appear in $\rho_1 \otimes \rho_{n/2-1}$.
\end{enumerate}
\end{proposition}
\begin{proof}
We use character inner products. For 2D irreps of $D_n$, the characters satisfy $\chi_{\rho_k}(a^l x^m) = 2\cos(\omega_k l)$ if $m = 0$ and $0$ if $m = 1$ (see Section~\ref{sec:dn_irreps}). The tensor product decomposition $\psi_1 \otimes \psi_2 \supseteq \psi$ is detected by the multiplicity
\[
\langle \chi_{\psi_1 \otimes \psi_2}, \chi_\psi \rangle
= \frac{1}{|G|} \sum_{g \in G} \chi_{\psi_1}(g)\, \chi_{\psi_2}(g)\, \overline{\chi_\psi(g)}.
\]

For two 2D irreps $\rho_i, \rho_j$:
\begin{align*}
\langle \chi_{\rho_i \otimes \rho_j}, \chi_{\rho_k} \rangle
&= \frac{1}{2n} \sum_{l=0}^{n-1} 4\cos(\omega_i l)\cos(\omega_j l) \cdot 2\cos(\omega_k l) \\
&= \frac{4}{2n} \sum_{l=0}^{n-1} \bigl[\cos(\omega_{i+j} l) + \cos(\omega_{|i-j|} l)\bigr] \cos(\omega_k l).
\end{align*}
Using $\sum_{l=0}^{n-1} \cos(\omega_m l) = n$ if $n \mid m$ and $0$ otherwise, this is nonzero if and only if $k = i + j$ or $k = |i - j|$ (with indices taken modulo $n$ and folded into the range $\{1, \ldots, \lfloor(n-1)/2\rfloor\}$). That is, the standard tensor-product rule for 2D dihedral irreps is
\[
\rho_i \otimes \rho_j \;\cong\; \rho_{i+j} \oplus \rho_{|i-j|}
\]
with boundary conventions as above. In particular, $\rho_k$ appears in $\rho_1 \otimes \rho_{k-1}$.

For the 1D irreps, $\rho_1 \otimes \rho_1$ decomposes as $\rho_0 \oplus \rho_{01} \oplus \rho_2$ (when $2 < n/2$), which is verified by $\langle \chi_{\rho_1 \otimes \rho_1}, \chi_{\rho_{01}} \rangle = 1$. For $n$ even, $\rho_1 \otimes \rho_{n/2-1}$ has $1 + (n/2 - 1) = n/2$ as the sum index, hitting the boundary index $k = n/2$, at which the corresponding 2D construction becomes reducible and decomposes as $\rho_{02} \oplus \rho_{03}$.
\end{proof}

\subsection{Constructive recovery procedure for $D_n$}

The following procedure recovers $\FF(\theta_x)$ from $\beta(x)$ for $G = D_n$, under the non-degeneracy assumption $\det(\FF(\theta_x)_\psi) \neq 0$ for all $\psi \in \Ghat$.

\textbf{Step 1.} Recover $\FF(\theta_x)_{\rho_0}$ from $\beta(x)_{\rho_0,\rho_0}$ using Lemma~\ref{lem:recover_rho0} and \eqref{eq:recover_rho0}.

\textbf{Step 2.} Recover $\FF(\theta_x)_{\rho_1}$ (up to unitary $W$) from $\beta(x)_{\rho_0, \rho_1}$ using Lemma~\ref{lem:recover_rho1} and \eqref{eq:recover_rho1}.

\textbf{Step 3.} For $k = 1, \ldots, \lfloor(n-1)/2\rfloor$: use the bootstrap (Lemma~\ref{lem:bootstrap}) with $\beta(x)_{\rho_1, \rho_k}$ to recover the Fourier blocks of all irreps appearing in $\rho_1 \otimes \rho_k$. By Proposition~\ref{prop:generation}:
\begin{itemize}
    \item $\beta_{\rho_1, \rho_1}$ yields $\rho_2$ and the 1D irrep $\rho_{01}$.
    \item $\beta_{\rho_1, \rho_k}$ for $k = 2, \ldots, \lfloor(n-1)/2\rfloor - 1$ yields $\rho_{k+1}$.
    \item $\beta_{\rho_1, \rho_{n/2-1}}$ (for $n$ even) yields $\rho_{02}$ and $\rho_{03}$.
\end{itemize}

\textbf{Step 4.} All Fourier blocks $\FF(\theta_x)_\psi$ are now known. Apply Fourier inversion (Proposition~\ref{prop:inversion}) to recover $\theta_x$ up to the $G$-action encoded in $W$.

Completeness of this procedure is guaranteed by Theorem~\ref{thm:completeness} (via Kakarala). The procedure provides an explicit algorithm; by Theorem~\ref{thm:completeness}, the residual unitary ambiguity $W$ from Step~2 corresponds globally to a single group translate.

\section{Selectivity}\label{sec:selectivity}

\subsection{The selective bispectrum}

The recovery procedure of the previous section uses only a small subset of bispectral coefficients. We formalize this.

\begin{definition}[Selective $G$-bispectrum]\label{def:selective}
A \textit{selective $G$-bispectrum} is a subset $\beta_{\mathrm{sel}}(\theta) \subset \{\beta(\theta)_{\psi_1,\psi_2}\}_{(\psi_1,\psi_2) \in \Ghat \times \Ghat}$ that is sufficient for reconstruction, i.e., from which $\FF(\theta)$ can be recovered (up to $G$-action) under the non-degeneracy assumption.
\end{definition}

\subsection{Selectivity for dihedral groups}

\begin{theorem}[Selectivity for $D_n$]\label{thm:selectivity}
For $G = D_n$ with $n \geq 4$ even, the per-pixel selective $G$-bispectrum requires at most $\lfloor(n-1)/2\rfloor + 2$ bispectral blocks under the assumption $\det(\FF(\theta_x)_\psi) \neq 0$ for all $\psi \in \Ghat$. The total number of matrix entries is $1 + 4 + 16\lfloor(n-1)/2\rfloor$.
\end{theorem}
\begin{proof}
The selective set consists of:
\begin{itemize}
    \item $\beta_{\rho_0,\rho_0}$: a $d_{\rho_0}^2 \times d_{\rho_0}^2 = 1 \times 1$ matrix (1 entry).
    \item $\beta_{\rho_0,\rho_1}$: a $d_{\rho_0}d_{\rho_1} \times d_{\rho_0}d_{\rho_1} = 2 \times 2$ matrix (4 entries).
    \item $\beta_{\rho_1,\rho_k}$ for $k = 1, \ldots, \lfloor(n-1)/2\rfloor$: each a $d_{\rho_1}d_{\rho_k} \times d_{\rho_1}d_{\rho_k} = 4 \times 4$ matrix (16 entries each). This is $\lfloor(n-1)/2\rfloor$ coefficients.
\end{itemize}
Total bispectral blocks: $2 + \lfloor(n-1)/2\rfloor$.

Total entries: $1 + 4 + 16 \cdot \lfloor(n-1)/2\rfloor$.
\end{proof}

\begin{remark}
We state the count for even $n$; the odd case is analogous and slightly smaller, since the 1D irreps $\rho_{02}, \rho_{03}$ do not exist and the selective set can be reduced accordingly. The stated bound remains valid for odd $n$.
\end{remark}

\subsection{Concrete counts}

The full bispectrum has $\sum_{(\psi_1,\psi_2)} (d_{\psi_1} d_{\psi_2})^2 = \bigl(\sum_\psi d_\psi^2\bigr)^2 = |G|^2$ entries.

\begin{table}[h]
\centering
\begin{tabular}{lrrrr}
\toprule
\textbf{Group} & $|G|$ & \textbf{Full bisp.\ entries} & \textbf{Sel.\ blocks} & \textbf{Sel.\ entries} \\
\midrule
$D_4$ & 8 & 64 & 3 & 21 \\
$D_8$ & 16 & 256 & 5 & 53 \\
$D_{16}$ & 32 & 1024 & 9 & 117 \\
$D_{32}$ & 64 & 4096 & 17 & 245 \\
\bottomrule
\end{tabular}
\caption{Comparison of full and selective per-pixel bispectrum sizes for dihedral groups. Entries are complex matrix entries.}
\label{tab:counts}
\end{table}

\begin{corollary}\label{cor:scaling}
The selective per-pixel $G$-bispectrum uses $O(|G|)$ entries per pixel, versus $O(|G|^2)$ for the full bispectrum.
\end{corollary}

\section{Discussion}

\subsection{Complete data flow}

The construction yields the following pipeline:
\[
f : \Z^2 \to \C^{K_{\mathrm{in}}}
\;\xrightarrow{\text{steerable conv}}\;
\Theta : \Z^2 \to \C[G]
\;\xrightarrow{\text{per-pixel } G\text{-FT}}\;
\widehat{\Theta}(x, \psi)
\;\xrightarrow{\text{sel.\ bisp.}}\;
\beta : \Z^2 \to \C^{O(|G|)}.
\]

\begin{itemize}
    \item $\beta$ is \textbf{$G$-invariant}: Theorem~\ref{thm:perpixel_bsp_inv}.
    \item $\beta$ is \textbf{translation-equivariant}: inherited from the convolutional structure.
    \item $\beta$ is a \textbf{generically complete} invariant of the fiber signal $\theta_x$: Theorem~\ref{thm:completeness}.
    \item $\beta$ uses $O(|G|)$ entries per pixel: Corollary~\ref{cor:scaling}.
\end{itemize}

Completeness for the original signal $f$ requires that the steerable convolution layer produce a sufficiently rich fiber signal $\theta_x$; this depends on the learned filters and is not addressed by the bispectral analysis alone.

\subsection{Relationship to the full-group approach}

An alternative approach defines the Fourier transform over the full group $\Gamma = \Z^2 \ltimes G$, integrating over all translations and all elements of $G$ simultaneously. This produces a global descriptor without a spatial index: the resulting bispectrum is invariant to the full group $\Gamma$ (translations and $G$ simultaneously), but it is not a spatial feature map and cannot be composed as a layer in a deep convolutional network.

Our approach decomposes the problem: spatial translation equivariance is handled by the convolutional structure, and $G$-invariance is handled per-pixel by the bispectrum. This yields a composable $G$-invariant layer that serves as a drop-in replacement for norm or gated nonlinearities in steerable architectures.

\subsection{Open questions}

The construction requires that the steerable CNN framework (e.g., \texttt{escnn}) supports setting the output fiber type to the regular representation $\rho_{\mathrm{reg}}$ of $G$. In \texttt{escnn}, this corresponds to setting the output type to $A_1 \oplus A_2 \oplus B_1 \oplus B_2 \oplus 2E$ for $G = D_4$ (total fiber dimension 8). Verifying that this is compatible with the intertwiner basis computation and the learned filter parameterization is an implementation question that remains to be validated.

\begin{thebibliography}{10}

\bibitem{kakarala1992}
R.~Kakarala.
\newblock Triple correlation on groups.
\newblock PhD thesis, University of California, Irvine, 1992.

\bibitem{kakarala2009}
R.~Kakarala.
\newblock Completeness of bispectrum on compact groups.
\newblock \textit{IEEE Transactions on Information Theory}, 2009.

\bibitem{mataigne2024}
S.~Mataigne, J.~Mathe, S.~Sanborn, C.~Shewmake, and N.~Miolane.
\newblock The selective {$G$}-bispectrum and its inversion: Applications to {$G$}-invariant networks.
\newblock In \textit{NeurIPS}, 2024.

\bibitem{weiler2019}
M.~Weiler and G.~Cesa.
\newblock General {E(2)}-equivariant steerable {CNNs}.
\newblock In \textit{NeurIPS}, 2019.

\bibitem{kondor2008}
R.~Kondor.
\newblock Group theoretical methods in machine learning.
\newblock PhD thesis, Columbia University, 2008.

\bibitem{sanborn2023bispectral}
S.~Sanborn, C.~Shewmake, and N.~Miolane.
\newblock Bispectral neural networks.
\newblock In \textit{ICLR}, 2023.

\end{thebibliography}

\end{document}
